\chapter{Archaeological Stratigraphy: The IUP Horizon}

\section{The Material Proxy for PED Radiation}
Historical linguistics cannot exist in a vacuum. To validate a 40,000-year-old reconstruction, we must identify a corresponding material radiation in the archaeological record. We identify the \textbf{Initial Upper Paleolithic (IUP)} lithic complex (c. 45,000–35,000 BP) as the definitive material proxy for the Proto-Eurasian-Dravidian radiation. This period marks a fundamental shift from Middle Paleolithic flake technologies to the sophisticated, modular blade-and-point assemblages that define the Upper Paleolithic.

\section{The Altai-Pamir Interaction Sphere}
The 'Pamir-Altai Corridor' served as the geographic hub where the PED mother tongue achieved its structural maturity. We identify two primary anchor sites:
\begin{enumerate}
    \item \textbf{Kara-Bom (Altai Mountains):} Excavations here document a continuous transition from Levallois-like techniques to specialized blade production c. 43,000 BP. This site represents the 'Phonological Laboratory' of the ancestral speakers.
    \item \textbf{Obi-Rakhmat (Tian Shan):} A critical site documenting the intersection of Siberian and Western Asian lithic traditions, providing the technological diversity that matches the PED morphological complexity.
\end{enumerate}

\subsection{Micro-Blade Dispersal: The Triple Radiation}
The stratigraphic evidence shows a rapid dispersal of IUP technologies along three primary axes, matching the linguistic divergence:
\begin{itemize}
    \item \textbf{The Eastern Axis (Sino-Tibetan):} Dispersal into the Tibetan Plateau and the Yellow River basin (e.g., Shuidonggou site, c. 38,000 BP).
    \item \textbf{The Southern Axis (Dravidian):} Dispersal toward the Indian subcontinent via the Pamir-Hindu Kush interface (e.g., Jwalapuram evidence).
    \item \textbf{The Western Axis (Indo-European):} Dispersal into the Pontic-Caspian steppe and Eastern Europe (e.g., Bacho Kiro cave, c. 45,000 BP).
\end{itemize}

\section{The Cognitive Revolution and Syntactic Stability}
We argue that the 'Cognitive Revolution' associated with the IUP—characterized by the first systematic use of bone tools, beads, and pigments—is the neurological substrate that allowed for the crystallization of the PED Active-Stative 'operating system'. The modularity of the lithic toolsets (blades stacked into composite tools) mirrored the modularity of the 'brick-stack' syntax.

\section{Conclusion: The Stratigraphic Signal}
The PED signal is not a ghost; it is a stratigraphic reality. The convergence of the Glottalic Shift and the IUP micro-blade expansion provides the multi-disciplinary proof required for scientific acceptance. The stones of Kara-Bom provide the material witness to the voice of the ancestors.
